\begin{figure}[H]
\centering
\begin{tikzpicture}[
    scale=0.8,
    vertex/.style={circle, fill=black, inner sep=1.6pt},
    edge/.style={line width=0.55pt},
    arr/.style={-{Latex[length=2mm]}, line width=0.55pt}
]

% =========================
% Graph 1
% =========================
\begin{scope}[xshift=0cm, yshift=0cm]

\node[vertex] (a11) at (0,3) {};
\node[vertex] (a21) at (1.2,3) {};
\node[vertex] (a31) at (2.4,3) {};

\node[vertex] (a12) at (0,2) {};
\node[vertex] (a22) at (1.2,2) {};
\node[vertex] (a32) at (2.4,2) {};

\node[vertex] (a13) at (0,1) {};
\node (a23) at (1.2,1) {$\times$};
\node[vertex] (a33) at (2.4,1) {};

\node[vertex] (a14) at (0,0) {};
\node[vertex] (a24) at (1.2,0) {};
\node[vertex] (a34) at (2.4,0) {};

\draw[edge] (a11)--(a21)--(a31);
\draw[edge] (a12)--(a22)--(a32);
\draw[edge] (a13)--(a23)--(a33);
\draw[edge] (a14)--(a24)--(a34);

\draw[edge] (a11)--(a12)--(a13)--(a14);
\draw[edge] (a21)--(a22)--(a23)--(a24);
\draw[edge] (a31)--(a32)--(a33)--(a34);

% \node at (1.45,1.02) {\(\star\)};

\end{scope}

\draw[arr] (3.1,1.5) -- (4.2,1.5);

% =========================
% Graph 2
% =========================
\begin{scope}[xshift=4.8cm, yshift=0.05cm]

\node[vertex] (b11) at (0,3) {};
\node[vertex] (b21) at (1.15,3) {};
\node[vertex] (b31) at (2.3,3) {};

\node[vertex] (b12) at (0,2) {};
\node[vertex] (b22) at (1.15,2) {};
\node[vertex] (b32) at (2.3,2) {};

\node[vertex] (b13) at (0,1) {};
\node[vertex] (b33) at (2.3,1) {};

\node[vertex] (b14) at (0,0) {};
\node[vertex] (b24) at (1.15,0) {};
\node[vertex] (b34) at (2.3,0) {};

\draw[edge] (b11)--(b21)--(b31);
\draw[edge] (b12)--(b22)--(b32);
\draw[edge] (b14)--(b24)--(b34);

\draw[edge] (b11)--(b12)--(b13)--(b14);
\draw[edge] (b21)--(b22);
\draw[edge] (b31)--(b32)--(b33)--(b34);

\node at (-0.3,0) {\(\star\)};
\node at (2.6,0) {\(\star\)};

\end{scope}

\draw[arr] (7.75,1.5) -- (8.85,1.5);

% =========================
% Graph 3
% =========================
\begin{scope}[xshift=9.4cm, yshift=0.05cm]

\node[vertex] (b11) at (0,3) {};
\node[vertex] (b21) at (1.15,3) {};
\node[vertex] (b31) at (2.3,3) {};

\node[vertex] (b12) at (0,2) {};
\node[vertex] (b22) at (1.15,2) {};
\node[vertex] (b32) at (2.3,2) {};

\node[vertex] (b13) at (0,1) {};
\node[vertex] (b33) at (2.3,1) {};

\node (b14) at (0,0) {$\times$};
\node[vertex] (b24) at (1.15,0) {};
\node (b34) at (2.3,0) {$\times$};

\draw[edge] (b11)--(b21)--(b31);
\draw[edge] (b12)--(b22)--(b32);
\draw[edge] (b14)--(b24)--(b34);

\draw[edge] (b11)--(b12)--(b13)--(b14);
\draw[edge] (b21)--(b22);
\draw[edge] (b31)--(b32)--(b33)--(b34);

\draw[edge] (b13) -- (b24);
\draw[edge] (b24) -- (b33);

\node at (-0.3,1) {\(\star\)};
\node at (2.6,1) {\(\star\)};

\end{scope}

\draw[arr] (11.9,1.5) -- (13.0,1.5);

% =========================
% Graph 4
% =========================
\begin{scope}[xshift=13.5cm, yshift=0.15cm]

\node[vertex] (b11) at (0,3) {};
\node[vertex] (b21) at (1.15,3) {};
\node[vertex] (b31) at (2.3,3) {};

\node[vertex] (b12) at (0,2) {};
\node[vertex] (b22) at (1.15,2) {};
\node[vertex] (b32) at (2.3,2) {};

\node[] (b13) at (0,1) {$\times$};
\node[] (b33) at (2.3,1) {$\times$};

\node[vertex] (b23) at (1.15,1) {};

\draw[edge] (b11)--(b21)--(b31);
\draw[edge] (b12)--(b22)--(b32);

\draw[edge] (b11)--(b12)--(b13);
\draw[edge] (b21)--(b22);
\draw[edge] (b31)--(b32)--(b33);

\draw[edge] (b13) -- (b23);
\draw[edge] (b23) -- (b33);

\draw[edge] (b12) -- (b23) -- (b32);

% \node at (-0.3,1) {\(\star\)};
% \node at (2.6,1) {\(\star\)};

\end{scope}

\draw[arr] (16.2,1.5) -- (17.3,1.5);

% =========================
% Graph 5
% =========================
\begin{scope}[xshift=17.8cm, yshift=0.2cm]

\node[vertex] (b11) at (0,3) {};
\node[vertex] (b21) at (1.15,3) {};
\node[vertex] (b31) at (2.3,3) {};

\node[vertex] (b12) at (0,2) {};
\node[vertex] (b22) at (1.15,2) {};
\node[vertex] (b32) at (2.3,2) {};

\node[vertex] (b23) at (1.15,1) {};

\draw[edge] (b11)--(b21)--(b31);
\draw[edge] (b12)--(b22)--(b32);

\draw[edge] (b11)--(b12);
\draw[edge] (b21)--(b22);
\draw[edge] (b31)--(b32);

\draw[edge] (b12) -- (b23) -- (b32);

\end{scope}
\end{tikzpicture}
\caption{Vertices marked by \(\star\) are locally complemented, and vertices marked by \(\times\) are deleted. This shows that $\Grid_{3 \times 4}$ has a forbidden minor of circle graphs as a vertex minor.}
\label{fig:grid-has-forbidden-vertex-minor}
\end{figure}